% Generated draft: requires mathematical and source review.
\documentclass[11pt]{article}
\usepackage[margin=25mm]{geometry}
\usepackage{fontspec}
\setmainfont{Latin Modern Roman}
\usepackage{amsmath,amssymb,mathtools,unicode-math}
\setmathfont{Latin Modern Math}
\usepackage{enumitem,needspace,xurl,hyperref}
\setlength{\parindent}{0pt}
\setlength{\parskip}{5pt}
\setlength{\emergencystretch}{3em}
\newcommand{\N}{\mathbb{N}}
\newcommand{\Z}{\mathbb{Z}}
\newcommand{\Q}{\mathbb{Q}}
\newcommand{\R}{\mathbb{R}}
\newcommand{\C}{\mathbb{C}}
\newcommand{\F}{\mathbb{F}}
\newcommand{\A}{\mathbb{A}}
\newcommand{\PP}{\mathbb{P}}
\newcommand{\E}{\mathbb{E}}
\newcommand{\eps}{\varepsilon}
\newcommand{\dd}{\,\mathrm{d}}
\begin{document}
\section*{UnsolvedMath ID: 30006561}
\textbf{Source problem code:} catalogue entry
\subsection{Definitions and mathematical statement}
The supplied record concerns a quantity denoted by \(C_d(w)\), described as an extremal word-occurrence constant for grids in dimension \(d\), with word parameter \(w\). It does not define the admissible grids, alphabet, occurrences, extremal normalization, limiting procedure, or periodicity. Consequently, \(C_d(w)\) cannot be given a self-contained mathematical definition from the supplied evidence.

Subject to the missing definitions, the recorded questions are:
\begin{enumerate}
\item For every dimension \(d\) and every word \(w\) admitted by the intended framework, does there exist a periodic grid construction attaining \(C_d(w)\) asymptotically?
\item Determine \(C_d(ABB)\) for every dimension \(d>2\).
\item Determine \(C_2(ABBB)\).
\end{enumerate}
Here \(ABB\) and \(ABBB\) denote the word patterns specified in the record. The alphabet and interpretation of these patterns must be recovered from the cited sources. In the first question, the periodic construction may depend on both \(d\) and \(w\). The meaning of asymptotic attainment, including the finite regions and boundary conventions used to evaluate a periodic construction, remains unspecified. These are source-level questions requiring completion of their definitions before a precise mathematical formulation is possible.
\subsection{Short English statement}
The record asks whether periodic constructions achieve every extremal grid word-occurrence constant asymptotically, and asks for \(C_d(ABB)\) when \(d>2\) and for \(C_2(ABBB)\). The definitions needed to make these questions precise are absent. This is a statement-only draft requiring human mathematical review.
\subsection{Sources}
\begin{itemize}
\item Supplied local problem registry record: permanent UnsolvedMath ID \(30006561\), source problem code \texttt{catalogue entry}, titled ``Periodic Constructions for Colored-Lattice Packing Constants.'' The record gives proposed year 2026 and names no proposer.
\item Supplied report citation: \emph{Combinatorics} (2026), Oberwolfach Reports, DOI 10.4171/owr/2026/1, pp.\ 5--79. Literature-source label: ``Oberwolfach Report 1/2026.'' \url{https://doi.org/10.4171/owr/2026/1}.
\item Supplied literature reference: Halberstam--Schildkraut, \emph{Words with repeated letters in a grid}. \url{https://arxiv.org/abs/2511.19678}.
\item Supplied literature assessment, dated exactly \texttt{2026-08-22}: the record labels the problem open, reports that only selected word and dimension cases are known, and states that general periodic attainment and the constants \(C_d(ABB)\) and \(C_2(ABBB)\) remain open. This is supplied dated curation derived from the Oberwolfach material, not an independently verified literature assessment; its status and cited claims require independent verification. The explicit question about \(C_d(ABB)\) is restricted to \(d>2\).
\item The registry labels its statement \texttt{corrected\_verified} and reports that the official report was checked, the question isolated from surrounding exposition, and its notation repaired for self-contained publication. These are supplied verification claims; no source visit or independent verification is asserted here. The supplied excerpt nevertheless omits the definitions needed for a self-contained formulation.
\end{itemize}
\end{document}
